\documentclass[11pt]{article}

\usepackage[T1]{fontenc}
\usepackage[utf8]{inputenc}
\usepackage{lmodern}
\usepackage[margin=1.05in]{geometry}
\usepackage{amsmath,amssymb,amsthm,mathtools}
\usepackage{aliascnt}
\usepackage{booktabs}
\usepackage{array}
\usepackage{enumitem}
\usepackage{float}
\usepackage{microtype}
\usepackage{xcolor}
\usepackage{tikz}
\usetikzlibrary{arrows.meta,calc,positioning}
\usepackage{hyperref}
\usepackage[nameinlink,noabbrev]{cleveref}

\definecolor{sourceblue}{HTML}{2563EB}
\definecolor{targetred}{HTML}{DC2626}
\definecolor{softgray}{HTML}{64748B}
\definecolor{darkink}{HTML}{172033}

\hypersetup{
  colorlinks=true,
  linkcolor=sourceblue,
  citecolor=sourceblue,
  urlcolor=sourceblue,
  pdftitle={One Edge, Unbounded Instability: Hilbert Decomposition Signed Measures in Multiparameter Persistence},
  pdfauthor={Anton Künzi},
  pdfsubject={Multiparameter persistent homology; Hilbert decomposition signed measures; formal verification},
  pdfkeywords={multiparameter persistence, Hilbert function, signed measure, Wasserstein distance, Lean}
}

\newtheorem{theorem}{Theorem}[section]
\newaliascnt{proposition}{theorem}
\newtheorem{proposition}[proposition]{Proposition}
\aliascntresetthe{proposition}
\newaliascnt{lemma}{theorem}
\newtheorem{lemma}[lemma]{Lemma}
\aliascntresetthe{lemma}
\newaliascnt{corollary}{theorem}
\newtheorem{corollary}[corollary]{Corollary}
\aliascntresetthe{corollary}
\theoremstyle{definition}
\newaliascnt{definition}{theorem}
\newtheorem{definition}[definition]{Definition}
\aliascntresetthe{definition}
\newaliascnt{remark}{theorem}
\newtheorem{remark}[remark]{Remark}
\aliascntresetthe{remark}

\newcommand{\R}{\mathbb{R}}
\newcommand{\Z}{\mathbb{Z}}
\newcommand{\N}{\mathbb{N}}
\newcommand{\1}{\mathbf{1}}
\newcommand{\KR}{\mathrm{KR}}
\newcommand{\Hil}{\mathrm{Hil}}
\newcommand{\dI}{d_{\mathcal I}^{1}}
\newcommand{\supp}{\operatorname{supp}}
\newcommand{\join}{\mathbin{\vee}}
\newcommand{\doi}[1]{\href{https://doi.org/#1}{\nolinkurl{doi:#1}}}
\newcommand{\fm}{f_m}
\newcommand{\gm}{g_m}
\newcommand{\Sm}{S_m}

\newcommand{\papertitlelead}{One Edge, Unbounded Instability:}
\newcommand{\papertitlesub}{Hilbert Decomposition Signed Measures in Multiparameter Persistence}

\newcommand{\makepapermodernfrontmatter}{%
  \vspace*{-2.2em}
  \noindent
  \begin{minipage}{\textwidth}
    \sffamily
    \textcolor{softgray}{\small\bfseries PREPRINT}
    \hfill\textcolor{softgray}{\small 30 July 2026}\par
    \vspace{0.7em}
    \textcolor{sourceblue}{\rule{\textwidth}{1.2pt}}\par
    \vspace{0.75em}
    {\fontsize{23}{27}\selectfont\bfseries\color{darkink}\papertitlelead\par}
    \vspace{0.2em}
    {\Large\color{darkink}\papertitlesub\par}
    \vspace{0.85em}
    {\large\bfseries\color{darkink}Anton Künzi\par}
    \vspace{0.35em}
    {\small\color{darkink}Independent researcher, prev.\ ETH Zurich\par}
    {\small
     ORCID: \href{https://orcid.org/0009-0002-1301-4631}{0009-0002-1301-4631}
     \enspace\textcolor{softgray}{/}\enspace
     \href{mailto:anton.kuenzi@alumni.ethz.ch}{anton.kuenzi@alumni.ethz.ch}\par}
    \vspace{0.6em}
    \textcolor{softgray!45}{\rule{\textwidth}{0.5pt}}
  \end{minipage}
  \vspace{0.6em}
}

\newcommand{\makepapertraditionalfrontmatter}{%
  \vspace*{-1.5em}
  \begin{center}
    {\LARGE\bfseries\papertitlelead\par}
    \vspace{0.25em}
    {\Large\papertitlesub\par}
    \vspace{1.05em}
    \paperauthorblock
    \vspace{0.55em}
    {\small\textit{Preprint}
     \enspace\textperiodcentered\enspace 30 July 2026\par}
  \end{center}
  \paperauthornote
  \vspace{0.8em}
}

\ifdefined\paperoptiononea
  \newcommand{\paperauthorblock}{%
    {\large Anton Künzi\par}
    \vspace{0.35em}
    {\small Independent researcher, prev.\ ETH Zurich\par}
    {\small
     ORCID: \href{https://orcid.org/0009-0002-1301-4631}{0009-0002-1301-4631}
     \enspace\textperiodcentered\enspace
     \href{mailto:anton.kuenzi@alumni.ethz.ch}{anton.kuenzi@alumni.ethz.ch}\par}
  }
  \newcommand{\paperauthornote}{}
  \newcommand{\makepaperfrontmatter}{\makepapertraditionalfrontmatter}
\else
  \ifdefined\paperoptiononeb
    \newcommand{\paperauthorblock}{{\large Anton Künzi\footnotemark\par}}
    \newcommand{\paperauthornote}{%
      \footnotetext{Independent researcher, prev.\ ETH Zurich. ORCID:
      \href{https://orcid.org/0009-0002-1301-4631}{0009-0002-1301-4631}.
      \href{mailto:anton.kuenzi@alumni.ethz.ch}{anton.kuenzi@alumni.ethz.ch}.}
    }
    \newcommand{\makepaperfrontmatter}{\makepapertraditionalfrontmatter}
  \else
    \newcommand{\makepaperfrontmatter}{\makepapermodernfrontmatter}
  \fi
\fi

\begin{document}
\makepaperfrontmatter

\begin{abstract}
We prove that Hilbert decomposition signed measures, a computable summary of
multiparameter persistent homology, are not uniformly stable over all finite
filtrations once there are at least three parameters.  For every $n\geq3$ and
$m\geq1$, we construct two filtrations of the same finite graph that differ
only in the grade of one edge, yet
\[
  \lVert f_m-g_m\rVert_1=1,\qquad
  \bigl\lVert\mu_{H_0(f_m)}-\mu_{H_0(g_m)}\bigr\rVert_{\KR,1}=2m.
\]
The output distance is therefore unbounded at fixed input distance one.  This
settles an open question of Loiseaux et al.~\cite{LoiseauxEtAl2023} and gives a
\hyperref[cor:phase]{sharp dimensional cutoff}: the dimension-only estimate
holds for one and two parameters, but no finite dimension-only constant exists
for $n\geq3$.
Rescaling puts every grade in the unit cube with input distance $1/m$ and
output distance two, ruling out a graph-independent modulus of continuity.
The four-vertex diamond graph is the smallest violation; after subdivision,
the same mechanism occurs in lower-star filtrations.  The corresponding
module-level bound also fails.  The counterexample family and its main
quantitative consequences are
\hyperref[app:lean]{formalized and kernel-checked in Lean}.
\end{abstract}

\section{Introduction}

Multiparameter persistent homology indexes a filtration by vectors in
$\R^n$.  The resulting modules rarely admit barcodes, so topological data
analysis uses computable summaries such as Hilbert functions, multigraded
Betti numbers, rank invariants, and signed decompositions
\cite{BotnanLesnick2023,BotnanOppermannOudot2025}.  A Hilbert decomposition
encodes the Hilbert function as a finite signed point measure.  These measures
support stable vectorizations for machine learning~\cite{LoiseauxEtAl2023}
and are implemented in \texttt{multipers}~\cite{LoiseauxSchreiber2024}.

Loiseaux et al.\ proved that, for a finite simplicial complex $S$ and
monotone maps $f,g:S\to\R^n$,
\begin{equation}\label{eq:known-positive}
  \bigl\lVert\mu_{H_i(f)}-\mu_{H_i(g)}\bigr\rVert_{\KR,1}
  \leq n\lVert f-g\rVert_1
  \qquad(n\in\{1,2\}),
\end{equation}
and explicitly left the extension to $n>2$ open
\cite[Theorem~1(1)]{LoiseauxEtAl2023}.  Oudot and Scoccola proved the
corresponding one- and two-parameter module-level estimate
\cite[Theorem~1.5]{OudotScoccola2024}.  We show that both extensions fail in
every parameter dimension $n\geq3$.  The failure already occurs in degree-zero
persistent homology of graphs when the grade of a single edge changes.

\begin{theorem}[Filtered-graph counterexample]\label{thm:main}
Fix a coefficient field.  For every $n\geq3$ and every integer $m\geq1$,
there are a finite simple graph $\Sm$ and monotone filtrations
$\fm,\gm:\Sm\to\R^n$ such that
\[
  \lVert\fm-\gm\rVert_1=1,\qquad
  \bigl\lVert\mu_{H_0(\fm)}-\mu_{H_0(\gm)}\bigr\rVert_{\KR,1}=2m.
\]
Only one edge has different grades under $\fm$ and $\gm$.
\end{theorem}

The four-vertex case is displayed in \cref{sec:diamond}.  The general
construction and the proof of \cref{thm:main}, including the exact transport
calculation and the extension from three to $n$ parameters, are given in
\cref{sec:family}.

\begin{corollary}[Exact dimensional phase transition]\label{cor:phase}
A finite constant depending only on parameter dimension can bound the
Hilbert signed-measure distance by filtration $\ell^1$-displacement for all
finite filtered complexes if and only if $n\leq2$.
\end{corollary}

Equation~\eqref{eq:known-positive} gives the positive half of
\cref{cor:phase}, and \cref{thm:main} gives the negative half.  In fact, the
failure is not specific to the proposed constant $n$: no finite constant
depending only on dimension can work.  After scalar rescaling,
\cref{cor:no-modulus} rules out even a graph-independent modulus of
continuity on the unit cube.  The graph in our construction grows with $m$.
This is compatible with the fixed-complex Lipschitz theorem of Scoccola et
al., whose constant may depend on the complex
\cite[Proposition~E.3]{ScoccolaEtAl2024}.

The smallest violation of the proposed three-parameter constant occurs at
$m=2$.  Its graph is the diamond $K_4\setminus\{\text{one edge}\}$, and its
exact output cost is $4>3$.  By \cref{prop:minimal}, no simple graph with
fewer vertices or fewer edges can violate the constant.

The module-level bound fails as well.  Let $M_m=H_0(f_m)$ and
$N_m=H_0(g_m)$.  Their cellular presentations have the same incidence
matrix; one column label moves by $\ell^1$-distance one.  The path-metric
presentation distance of Bjerkevik and Lesnick
\cite{BjerkevikLesnick2021} then gives the following result.

\begin{theorem}[Presentation-distance counterexample]\label{thm:p3}
For every $n\geq3$ and $m\geq1$, there are finitely presented
$n$-parameter persistence modules $M_m,N_m$ such that
\[
  \dI(M_m,N_m)\leq1,\qquad
  d^{\Hil}_{W,1}(M_m,N_m)=2m.
\]
No finite dimension-only constant therefore bounds Hilbert signed
$1$-Wasserstein distance by $\dI$ for fixed $n\geq3$.
\end{theorem}

We prove \cref{thm:p3} in \cref{sec:p3}.  The unit-cube and lower-star
variants, the minimality result, and the $H_1$/Euler calculation are proved in
\cref{sec:diamond,sec:consequences}.  Reproducibility information appears in
\cref{sec:artifact};
\hyperref[app:lean]{Appendix~\ref*{app:lean}} records the correspondence
between the paper statements and their Lean formalization.

\section{Hilbert signed measures and transport}\label{sec:background}

Let $K$ be a finite simplicial complex.  A filtration
$f:K\to\R^n$ is \emph{monotone} when
$f(\tau)\leq f(\sigma)$ coordinatewise whenever $\tau$ is a face of
$\sigma$.  Its sublevel complex at $x\in\R^n$ is
\[
  K^f_x=\{\sigma\in K:f(\sigma)\leq x\}.
\]
Fixing a coefficient field, the persistent homology module has Hilbert
function
\[
  h_{H_i(f)}(x)=\dim H_i(K^f_x).
\]
For two filtrations of the same finite complex, we use the simplexwise
$\ell^1$ norm
\[
  \lVert f-g\rVert_1
  =\sum_{\sigma\in K}\lVert f(\sigma)-g(\sigma)\rVert_1.
\]

For the finitely presentable modules considered here, the
\emph{Hilbert decomposition signed measure} $\mu_M$ is the unique finite
signed point measure satisfying
\begin{equation}\label{eq:hilbert-decomp}
  h_M(x)=\mu_M((-\infty,x])
  \quad\text{for all }x\in\R^n,
  \qquad
  (-\infty,x]=\{y:y\leq x\}.
\end{equation}
This is the Möbius inverse of the Hilbert function
\cite[Definition~4 and Appendix~A.4]{LoiseauxEtAl2023}.  We write
$\mu_{H_i(f)}$ for the measure of $H_i(f)$.

If a finite signed atomic measure $\Delta$ has Jordan decomposition
$\Delta=\Delta^+-\Delta^-$ with equal positive and negative total mass,
Loiseaux et al.\ define its Kantorovich-Rubinstein norm with $\ell^1$ ground
cost as follows~\cite[Definition~3]{LoiseauxEtAl2023}:
\begin{equation}\label{eq:kr}
  \lVert\Delta\rVert_{\KR,1}
  =\inf_{\pi\in\Pi(\Delta^+,\Delta^-)}
      \int_{\R^n\times\R^n}\lVert x-y\rVert_1\,d\pi(x,y).
\end{equation}
All measures below are finite sums of unit atoms, so \eqref{eq:kr} is a
finite transportation problem.  The lower bound uses only the fact that
distinct integer-lattice points have distance at least one.  Explicit
matchings give the upper bound.

The Euler decomposition signed measure
\cite[Definition~5]{LoiseauxEtAl2023} is
\[
  \mu_{\chi(f)}=\sum_i(-1)^i\mu_{H_i(f)}.
\]
Unlike the degree-specific Hilbert measures, it satisfies the
dimension-independent estimate
\begin{equation}\label{eq:euler-positive}
  \lVert\mu_{\chi(f)}-\mu_{\chi(g)}\rVert_{\KR,1}
  \leq \lVert f-g\rVert_1
\end{equation}
in every parameter dimension~\cite[Theorem~1(2)]{LoiseauxEtAl2023}.

\section{The four-vertex counterexample}\label{sec:diamond}

The graph $S_2$ has vertices $s,t,u_1,u_2$, central edge
$e=\{s,t\}$, and two two-edge paths
$s-a_j-u_j-b_j-t$.  Every vertex has grade $(0,0,0)$, the arm grades agree
under $f_2$ and $g_2$, and only $e$ moves:
\[
\begin{array}{c|ccccc}
 & a_1 & b_1 & a_2 & b_2 & e\\
\midrule
 f_2 & (0,1,0)&(0,0,2)&(0,2,0)&(0,0,1)&(2,0,0)\\
 g_2 & (0,1,0)&(0,0,2)&(0,2,0)&(0,0,1)&(1,0,0).
\end{array}
\]
The filtered diamond is shown in \cref{fig:diamond}.

\begin{figure}[t]
\centering
\begin{tikzpicture}[
  vertex/.style={circle,draw=darkink,fill=white,minimum size=7mm,inner sep=1pt},
  common/.style={draw=darkink,line width=0.9pt},
  moved/.style={draw=targetred,line width=2.2pt},
  lab/.style={font=\scriptsize,fill=white,inner sep=1.5pt}
]
  \node[vertex] (s) at (-2.8,0) {$s$};
  \node[vertex] (t) at ( 2.8,0) {$t$};
  \node[vertex] (u1) at (0,1.7) {$u_1$};
  \node[vertex] (u2) at (0,-1.7) {$u_2$};
  \draw[common] (s)--node[lab,above left] {$(0,1,0)$} (u1);
  \draw[common] (u1)--node[lab,above right] {$(0,0,2)$} (t);
  \draw[common] (s)--node[lab,below left] {$(0,2,0)$} (u2);
  \draw[common] (u2)--node[lab,below right] {$(0,0,1)$} (t);
  \draw[moved] (s)--node[lab,above] {$e:\ (2,0,0)\mapsto(1,0,0)$} (t);
  \node[font=\small,align=center,text=softgray] at (0,-2.55)
    {$\lVert f_2-g_2\rVert_1=1$ \qquad
     $\lVert\Delta_2\rVert_{\KR,1}=4$};
\end{tikzpicture}
\caption{The diamond graph is the minimal counterexample: one edge moves by
one, while the degree-zero Hilbert signed measure moves by four.}
\label{fig:diamond}
\end{figure}

The signed-measure difference $\Delta_2=\mu_{H_0(f_2)}-\mu_{H_0(g_2)}$
has positive atoms
\[
 (1,0,0),\ (1,2,2),\ (2,1,2),\ (2,2,1)
\]
and negative atoms
\[
 (1,1,2),\ (1,2,1),\ (2,0,0),\ (2,2,2).
\]
The supports are disjoint subsets of $\Z^3$, so every coupling costs at
least four.  Matching vertically between the layers $r=1$ and $r=2$ costs
exactly four.  Since only $e$ moves, this proves $4>3\cdot1$.

\begin{proposition}[Minimality of the diamond]\label{prop:minimal}
Among simple graphs with arbitrary monotone three-parameter filtrations, a
counterexample to the proposed constant $3$ requires at least four vertices
and at least five edges.  Hence $S_2$ is simultaneously vertex-minimal and
edge-minimal.
\end{proposition}

\begin{proof}
A simple graph with at most four edges has cycle rank at most one.  If it is
a forest, then $\mu_{H_1}=0$ and
$\mu_{H_0}=\mu_\chi$, so \eqref{eq:euler-positive} gives the stronger
constant $1$.

If the graph is unicyclic, let $C$ be its unique cycle.  Its $H_1$ Hilbert
function is the indicator of the upper orthant born at
\[
  b_f=\bigvee_{e\in C}f(e).
\]
Monotonicity makes the vertex grades irrelevant once the cycle edges are
present.  Hence $\mu_{H_1(f)}=\delta_{b_f}$, and similarly for $g$.  For each
coordinate, the maximum operation is $1$-Lipschitz, hence
\[
  \lVert b_f-b_g\rVert_1
  \leq\sum_{e\in C}\lVert f(e)-g(e)\rVert_1
  \leq\lVert f-g\rVert_1.
\]
Using $\mu_{H_0}=\mu_\chi+\mu_{H_1}$ and the triangle inequality gives
\[
  \lVert\mu_{H_0(f)}-\mu_{H_0(g)}\rVert_{\KR,1}
  \leq2\lVert f-g\rVert_1.
\]
No graph with at most four edges can therefore violate the constant $3$.
A simple graph with at most three vertices has at most three edges, which
also proves vertex-minimality.
\end{proof}

\section{The general family}\label{sec:family}

\subsection{Graph and filtrations}

For $m\geq1$, let $\Sm$ have vertices
\[
  s,t,u_1,\ldots,u_m
\]
and edges
\[
  e=\{s,t\},\qquad a_j=\{s,u_j\},\qquad b_j=\{u_j,t\}
  \quad(1\leq j\leq m).
\]
There are no higher-dimensional simplices.  Give every vertex grade
$(0,0,0)$ and set
\begin{align*}
  \fm(a_j)=\gm(a_j)&=(0,j,0),\\
  \fm(b_j)=\gm(b_j)&=(0,0,m+1-j),\\
  \fm(e)&=(2,0,0),&
  \gm(e)&=(1,0,0).
\end{align*}
All coordinates are nonnegative, so the filtrations are monotone.  Only $e$
moves, and therefore
\begin{equation}\label{eq:input-one}
  \lVert\fm-\gm\rVert_1=1.
\end{equation}
The construction and the staircase measure arising below are summarized in
\cref{fig:family}.

\begin{figure}[t]
\centering
\begin{tikzpicture}[
  vertex/.style={circle,draw=darkink,fill=white,minimum size=6.5mm,inner sep=1pt},
  common/.style={draw=darkink,line width=0.8pt},
  moved/.style={draw=targetred,line width=2pt},
  atomplus/.style={circle,fill=sourceblue,draw=sourceblue,minimum size=4.5pt,inner sep=0pt},
  atomminus/.style={rectangle,fill=targetred,draw=targetred,minimum size=4.5pt,inner sep=0pt},
  lab/.style={font=\scriptsize,fill=white,inner sep=1pt}
]
  \begin{scope}[xshift=-4.4cm]
    \node[vertex] (s) at (-2.0,0) {$s$};
    \node[vertex] (t) at ( 2.0,0) {$t$};
    \node[vertex] (u1) at (0,1.65) {$u_1$};
    \node[vertex] (u2) at (0,0.65) {$u_2$};
    \node at (0,-0.25) {$\vdots$};
    \node[vertex] (um) at (0,-1.35) {$u_m$};
    \draw[common] (s)--(u1)--(t);
    \draw[common] (s)--(u2)--(t);
    \draw[common] (s)--(um)--(t);
    \draw[moved] (s)--node[lab,below] {$e$} (t);
    \node[font=\small,text=softgray] at (0,-2.05) {$S_m$: $m$ parallel two-edge paths};
  \end{scope}
  \begin{scope}[xshift=3.6cm,yshift=-1.55cm,scale=0.72]
    \draw[-{Latex[length=2mm]},darkink] (-0.3,0)--(6.1,0) node[right] {$y$};
    \draw[-{Latex[length=2mm]},darkink] (0,-0.3)--(0,6.1) node[above] {$z$};
    \foreach \j in {1,...,5}{
      \pgfmathtruncatemacro{\z}{6-\j}
      \node[atomminus] at (\j,\z) {};
    }
    \foreach \j in {1,...,4}{
      \pgfmathtruncatemacro{\x}{\j+1}
      \pgfmathtruncatemacro{\z}{6-\j}
      \node[atomplus] at (\x,\z) {};
    }
    \node[atomplus] at (0,0) {};
    \draw[dashed,softgray] (1,5)--(2,5)--(2,4)--(3,4)--(3,3)--(4,3)--(4,2)--(5,2)--(5,1);
    \node[font=\scriptsize,text=sourceblue,anchor=west] at (0.25,5.75) {positive: $0,c_j$};
    \node[font=\scriptsize,text=targetred,anchor=west] at (3.25,5.75) {negative: $p_j$};
    \node[font=\small,text=softgray] at (3,-1.0) {staircase measure $\nu_m$ (shown for $m=5$)};
  \end{scope}
\end{tikzpicture}
\caption{Left: the graph family.  Right: the two-dimensional staircase
Möbius measure; the blue circle at the origin and the blue corner atoms
alternate with the red path atoms.}
\label{fig:family}
\end{figure}

\subsection{The component-count formula}

Write a parameter as $x=(r,y,z)$.  Outside the slab $1\leq r<2$, the two
sublevel graphs agree.  Inside the slab, the $g_m$-sublevel graph contains
$e$ while the $f_m$-sublevel graph does not.  Adding $e$ merges two
components exactly when no path $s-u_j-t$ is complete.  The $j$th path is
complete exactly when
\[
  y\geq j,\qquad z\geq m+1-j.
\]
Define
\[
q_m(y,z)=
\begin{cases}
1,&y,z\geq0\ \text{and no }j\text{ satisfies }
       y\geq j,\ z\geq m+1-j,\\
0,&\text{otherwise}.
\end{cases}
\]
Since ordinary $H_0$ is the free vector space on connected components, over
every coefficient field,
\begin{equation}\label{eq:hilbert-difference}
  D_m(r,y,z)
  :=\dim H_0(K^{f_m}_{(r,y,z)})
     -\dim H_0(K^{g_m}_{(r,y,z)})
  =\1_{[1,2)}(r)q_m(y,z).
\end{equation}

\subsection{Möbius inversion of the staircase}

For $1\leq j\leq m$, put
\[
  p_j=(j,m+1-j),
\]
and, for $1\leq j<m$, put
\[
  c_j=p_j\join p_{j+1}=(j+1,m+1-j).
\]
Define the signed point measure
\begin{equation}\label{eq:nu}
  \nu_m=\delta_{(0,0)}
        -\sum_{j=1}^{m}\delta_{p_j}
        +\sum_{j=1}^{m-1}\delta_{c_j}.
\end{equation}

\begin{lemma}[Staircase inversion]\label{lem:staircase}
For every $(y,z)\in\R^2$,
\[
  \nu_m((-\infty,y]\times(-\infty,z])=q_m(y,z).
\]
\end{lemma}

\begin{proof}
If $y<0$ or $z<0$, both sides vanish.  Otherwise let
\[
  A(y,z)=\{j:p_j\leq(y,z)\}.
\]
Because the first coordinates of the $p_j$ increase while their second
coordinates decrease, $A(y,z)$ is empty or a consecutive interval.  If it
has $a>0$ elements, exactly $a-1$ adjacent joins $c_j$ lie below $(y,z)$.
The cumulative value of \eqref{eq:nu} is then
$1-a+(a-1)=0$.  If $A(y,z)$ is empty, its value is $1$.
\end{proof}

Let $\eta=\delta_1-\delta_2$ on $\R$.  Its cumulative function is
$\1_{[1,2)}$.  Equations \eqref{eq:hilbert-decomp},
\eqref{eq:hilbert-difference}, and \cref{lem:staircase} give
\begin{equation}\label{eq:delta-tensor}
  \Delta_m
  :=\mu_{H_0(f_m)}-\mu_{H_0(g_m)}
  =\eta\otimes\nu_m.
\end{equation}

Write
\[
  \nu_m^+=\delta_{(0,0)}+\sum_{j=1}^{m-1}\delta_{c_j},
  \qquad
  \nu_m^-=\sum_{j=1}^{m}\delta_{p_j}.
\]
Both have mass $m$, and no atom cancels in
\begin{align}
  \Delta_m^+
    &=\delta_1\otimes\nu_m^+
      +\delta_2\otimes\nu_m^-,
      \label{eq:positive}\\
  \Delta_m^-
    &=\delta_1\otimes\nu_m^-
      +\delta_2\otimes\nu_m^+.
      \label{eq:negative}
\end{align}
The origin, the $p_j$, and the $c_j$ are distinct in every potentially
cancelling pair.  Hence \eqref{eq:positive} and \eqref{eq:negative} are the
Jordan parts, each with $2m$ unit atoms.

\subsection{Exact transport and zero-padding}

The two supports in \eqref{eq:positive} and \eqref{eq:negative} are disjoint
subsets of $\Z^3$.  Every distance from a source atom to a target atom is at
least one, so every coupling costs at least $2m$.  For the reverse bound, match
$(1,a)$ to $(2,a)$ for every atom $a$ of $\nu_m^+$, and match
$(2,b)$ to $(1,b)$ for every atom $b$ of $\nu_m^-$.  Each of the $2m$ units
moves by one, giving
\[
  \lVert\Delta_m\rVert_{\KR,1}=2m.
\]
Equation~\eqref{eq:input-one} now proves \cref{thm:main} for three
parameters.  For $n>3$, append zero coordinates to every grade and atom.
Monotonicity, sublevel components, and both $\ell^1$ costs are unchanged.

\section{Consequences and robustness}\label{sec:consequences}

\subsection{No graph-independent modulus on the unit cube}

For $m\geq2$, divide every grade by $m$.  Denote the resulting filtrations
by $\widetilde f_m,\widetilde g_m$.  All grades lie in $[0,1]^n$ and scalar
dilation pushes forward the Hilbert signed measures.  Since the ground
metric is homogeneous,
\begin{equation}\label{eq:unit-cube}
  \lVert\widetilde f_m-\widetilde g_m\rVert_1=\frac1m,
  \qquad
  \bigl\lVert
    \mu_{H_0(\widetilde f_m)}
    -\mu_{H_0(\widetilde g_m)}
  \bigr\rVert_{\KR,1}=2.
\end{equation}

\begin{corollary}\label{cor:no-modulus}
For every fixed $n\geq3$, the Hilbert signed-measure map on all finite
$n$-parameter filtered graphs has no graph-independent modulus of uniform
continuity, even with grades restricted to $[0,1]^n$.  In particular, no
bound
\[
  \lVert\mu_{H_0(f)}-\mu_{H_0(g)}\rVert_{\KR,1}
  \leq C\lVert f-g\rVert_1^\alpha
\]
holds uniformly for any $C<\infty$ and $\alpha>0$.
\end{corollary}

\begin{proof}
In \eqref{eq:unit-cube}, the input tends to zero while the output remains
two.  The Hölder conclusion follows by taking $m$ large enough that
$C/m^\alpha<2$.
\end{proof}

Although the graph grows with $m$, it remains simple, planar,
two-terminal series-parallel, and of treewidth two.  Its size is
\[
  |V(S_m)|=m+2,\qquad |E(S_m)|=2m+1,\qquad |S_m|=3m+3
\]
when vertices and edges are both counted as simplices.  Along this family,
any valid complex-dependent replacement constant must grow at least
linearly: the ratio is
$2(|V|-2)=\tfrac23(|S_m|-3)$.

\subsection{Lower-star filtrations}

The freedom to grade edges independently is not essential.

\begin{proposition}[Lower-star realization]\label{prop:lower-star}
For every $n\geq3$ and $m\geq1$, there are vertex functions on a subdivided
finite graph, denoted $v_m^f$ and $v_m^g$, whose induced lower-star
filtrations $f_m^{\mathrm{LS}}$ and $g_m^{\mathrm{LS}}$ satisfy
\[
  \lVert v_m^f-v_m^g\rVert_1=1,\qquad
  \lVert f_m^{\mathrm{LS}}-g_m^{\mathrm{LS}}\rVert_1=3,
  \qquad
  \bigl\lVert
    \mu_{H_0(f_m^{\mathrm{LS}})}
    -\mu_{H_0(g_m^{\mathrm{LS}})}
  \bigr\rVert_{\KR,1}=2m.
\]
After division by $m$, the vertex-input distance is $1/m$, the induced
simplexwise distance is $3/m$, and the output distance is $2$.
\end{proposition}

\begin{proof}
Subdivide every edge of $S_m$ once.  Give each original vertex grade zero
and give the midpoint of an old edge its old edge grade.  In the induced
lower-star filtration, the midpoint and both incident subdivided edges enter
at exactly the old edge grade.  Replacing an active old edge by this active
two-edge path preserves connected components at every parameter, so the
$H_0$ Hilbert functions and signed measures are unchanged.

Only the midpoint of the central edge has different vertex grades.  Its own
grade and the grades of its two incident subdivided edges each move by one,
giving total filtration displacement three.  Rescaling gives the last
claim.
\end{proof}

\subsection[Same-family H1 and Euler cancellation]
{Same-family \(H_1\) and Euler cancellation}

The dimension-independent Euler bound \eqref{eq:euler-positive} may seem
surprising next to \cref{thm:main}.  The same graph family explains the
difference.
Set
\[
  \sigma_m=\sum_{j=1}^m\delta_{p_j}
          -\sum_{j=1}^{m-1}\delta_{c_j},
  \qquad
  \nu_m=\delta_{(0,0)}-\sigma_m.
\]
Let $\Delta_i=\mu_{H_i(f_m)}-\mu_{H_i(g_m)}$ and
$\Delta_\chi=\Delta_0-\Delta_1$.

\begin{proposition}[Atomwise Euler cancellation]\label{prop:euler-cancel}
For the family in \cref{thm:main},
\[
  \Delta_0=\eta\otimes(\delta_{(0,0)}-\sigma_m),\qquad
  \Delta_1=-\eta\otimes\sigma_m,\qquad
  \Delta_\chi=\eta\otimes\delta_{(0,0)}.
\]
Moreover,
\[
  \lVert\Delta_0\rVert_{\KR,1}=2m,\qquad
  \lVert\Delta_1\rVert_{\KR,1}=2m-1,\qquad
  \lVert\Delta_\chi\rVert_{\KR,1}=1.
\]
\end{proposition}

\begin{proof}
For a sublevel graph,
\[
  \beta_1=|E|-|V|+\beta_0.
\]
Inside the slab, $g_m$ has exactly one more active edge than $f_m$ whenever
the common vertex grade is active.  Hence the $H_1$ Hilbert-function
difference is the $H_0$ difference with the upper-orthant indicator born at
$(1,0,0)$ and dying at $(2,0,0)$ removed.  Möbius inversion gives the three
displayed identities.

The $H_1$ Jordan parts have $m+(m-1)=2m-1$ atoms each: corner atoms on one
layer and path atoms on the other.  Their supports are disjoint integer
lattices and the vertical matching has unit cost per atom.  The Euler
difference has just the two atoms $(1,0,0)$ and $(2,0,0)$.
\end{proof}

In the alternating Euler sum, all staircase atoms cancel and only the moved
simplex remains.

\section{The module-level presentation distance}\label{sec:p3}

For a filtered graph, the cellular sequence
\[
  C_1(f_m)\xrightarrow{\partial_1}C_0(f_m)
  \longrightarrow H_0(f_m)\longrightarrow0
\]
is a finite labeled presentation.  Choose the same orientations and simplex
ordering for $f_m$ and $g_m$.  The resulting matrices are identical.  Their
row labels and arm-column labels agree, while the central column label moves
from $(2,0,0)$ to $(1,0,0)$.

Bjerkevik and Lesnick define the preliminary $1$-presentation distance
between modules by minimizing the total $\ell^1$ label displacement over
compatible same-matrix presentations.  Their presentation distance
$d_{\mathcal I}^{1}$ is the path metric induced by this preliminary
dissimilarity~\cite[Section~3.1 and Definition~3.2]{BjerkevikLesnick2021};
see also the complexity analysis of Bjerkevik and
Botnan~\cite{BjerkevikBotnan2025}.
This compatible pair has cost one, so the one-step path gives
\[
  \dI(H_0(f_m),H_0(g_m))\leq1.
\]
The pointwise cokernel of the incidence presentation is ordinary cellular
$H_0$.  Its Hilbert function is therefore the component-count function from
\cref{sec:family}, whose signed-Wasserstein value is $2m$.  This proves
\cref{thm:p3}.

After scaling every presentation label by $1/m$, we have
\[
  \dI(M_m,N_m)\leq\frac1m,\qquad
  d^{\Hil}_{W,1}(M_m,N_m)=2.
\]
The module invariant therefore has no dimension-only modulus of uniform
continuity.  The argument uses only the upper bound on $\dI$; equality is
neither needed nor claimed.

\section{Artifact and reproducibility}\label{sec:artifact}

The paper is accompanied by a kernel-checked Lean formalization and an
independent finite computation.  The accompanying repository contains the
Lean source, build instructions, and pinned dependencies: Lean
\texttt{v4.33.0-rc1} and Mathlib at the revision recorded in the Lake manifest.
Running \texttt{lake build} from the repository root checks the complete formal
development; the README gives prerequisites and additional audit commands.  A
NetworkX/NumPy/SciPy script independently reproduces the $H_0$ component
counts, finite differences, and optimal assignments for $m=1,\ldots,15$.  The
computation is not used as a premise in Lean.

\paragraph{Use of AI-assisted tools.}
AI-assisted tools were used for literature triage, Lean formalization,
and manuscript editing.  The proofs were developed over the course of two
months from conversations between the author and GPT-5.5 / 5.6 Sol.  The
author checked all mathematical statements and proof artifacts and takes
responsibility for the content.

\paragraph{Competing interests.}
The author declares no competing interests.

\paragraph{Data and code availability.}
All constructions are finite and specified in the paper.  Lean source,
build metadata, and secondary verification scripts are available from
\url{https://github.com/ZeterMordio/one-edge-instability} under the
Apache-2.0 license.  \hyperref[app:lean]{Appendix~\ref*{app:lean}} gives the
theorem-to-source map and the formal trust audit.
% RELEASE GATE: Replace this pre-release sentence with the exact GitHub commit
% and Zenodo software DOI after the public release has been archived.
The archived software
release will identify its exact commit and DOI.  This preprint and its source
are licensed under
\href{https://creativecommons.org/licenses/by/4.0/}{CC BY~4.0}.

\appendix
\section{Lean formalization certificate}\label{app:lean}

The formal development uses Lean~\cite{deMouraUllrich2021} and
Mathlib~\cite{Mathlib2020}.  The principal $H_0$ source file,
\texttt{P1.lean}, imports only \texttt{Mathlib}.  It belongs to the
\texttt{OneEdgeInstability} library; the module-level result is proved
separately in \texttt{P3.lean}.  The development introduces the following
definitions for this family:

\begin{itemize}[leftmargin=1.4em]
  \item finite vertex, edge, and simplex types for $S_m$;
  \item real sublevel graphs and ordinary cellular $H_0=C_0/\operatorname{im}\partial_1$;
  \item finite signed atomic measures, lower-orthant cumulative functions,
        and uniqueness of Möbius inversion;
  \item conversion to Mathlib signed measures and identification of their
        Jordan positive and negative parts;
  \item $\KR_1$ as an infimum over Mathlib measures on
        $\R^n\times\R^n$, with measurable-set marginal equations and
        \texttt{lintegral} cost;
  \item labeled cellular presentations and the finite-path definition
        of $d_{\mathcal I}^1$;
  \item cellular $H_1=\ker\partial_1$, field-independent rank, and the
        Euler-Poincaré bridge.
\end{itemize}

\begin{table}[H]
\centering
\footnotesize
\begin{tabular}{>{\raggedright\arraybackslash}p{0.39\textwidth}
                >{\raggedright\arraybackslash}p{0.51\textwidth}}
\toprule
Paper statement & Lean declaration\\
\midrule
\Cref{thm:main} &
\texttt{P1.theorem1\_H0\_counterexample\_family}\\
Failure of the constant-$n$ bound &
\texttt{P1.proposed\_constant\_n\_P1\_inequality\_false}\\
No finite dimension-only constant &
\texttt{P1.no\_finite\_dimension\_only\_P2\_constant}\\
Diamond example $4>3$ &
\texttt{P1.n3\_m2\_explicit\_counterexample}\\
Unit-cube rescaling \eqref{eq:unit-cube} &
\texttt{P1.unit\_cube\_H0\_counterexample\_family}\\
Arbitrarily small inputs &
\texttt{P1.arbitrarily\_close\_unit\_cube\_counterexamples}\\
\Cref{prop:euler-cancel} &
\texttt{P1.same\_family\_H1\_and\_Euler\_cancellation}\\
\Cref{thm:p3} &
\texttt{P3.p3\_H0\_presentation\_counterexample\_family}\\
\bottomrule
\end{tabular}
\caption{Public theorem-to-paper map.}
\label{tab:lean-map}
\end{table}

All Lean names in \cref{tab:lean-map} share the root namespace
\texttt{OneEdgeInstability}, which is omitted there.

The chain-level $H_0$ proof identifies the quotient by the span of oriented
edge boundaries with the free vector space on connected components.  For
$H_1$, the proof identifies the range of the boundary with that same
cellular span, applies rank-nullity, and constructs the two Hilbert measures
before subtracting them.  The transport lower bound follows directly from
the marginal equations: every product-measure coupling is supported on the
finite source-to-target product, where the disjoint integer lattices impose
cost at least one per unit of mass.  The upper bound is the vertical
matching, written as a finite sum of Dirac measures on the product.

The checked source contains no \texttt{sorry}, \texttt{admit}, user-defined
\texttt{axiom}, or \texttt{unsafe} declaration.  For the audited public
theorems, \texttt{\#print axioms} reports only Lean's standard principles
\texttt{propext}, \texttt{Classical.choice}, and \texttt{Quot.sound}.

The positive $n\leq2$ half of \cref{cor:phase} is cited rather than
reformalized.  The minimality proof in \cref{prop:minimal} and the
subdivision argument in \cref{prop:lower-star} are paper proofs, not Lean
theorems.  Lean checks the counterexample family, the signed-measure and
transport identifications, the rescaling, the module presentation theorem,
and \cref{prop:euler-cancel}.

\begin{thebibliography}{99}
\raggedright
\interlinepenalty=10000

\bibitem{LoiseauxEtAl2023}
D.~Loiseaux, L.~Scoccola, M.~Carrière, M.~B.~Botnan, and S.~Oudot.
\newblock Stable vectorization of multiparameter persistent homology using
signed barcodes as measures.
\newblock In \emph{Advances in Neural Information Processing Systems},
volume~36, pages 68316--68342, 2023.
\newblock \doi{10.52202/075280-2987}.
\newblock \url{https://proceedings.neurips.cc/paper_files/paper/2023/hash/d75c474bc01735929a1fab5d0de3b189-Abstract-Conference.html}.

\bibitem{OudotScoccola2024}
S.~Oudot and L.~Scoccola.
\newblock On the stability of multigraded Betti numbers and Hilbert
functions.
\newblock \emph{SIAM Journal on Applied Algebra and Geometry},
8(1):54--88, 2024.
\newblock \doi{10.1137/22M1489150}.

\bibitem{BjerkevikLesnick2021}
H.~B.~Bjerkevik and M.~Lesnick.
\newblock $\ell^p$-distances on multiparameter persistence modules.
\newblock arXiv:2106.13589v2, 2021.
\newblock \url{https://arxiv.org/abs/2106.13589}.

\bibitem{BotnanLesnick2023}
M.~B.~Botnan and M.~Lesnick.
\newblock An introduction to multiparameter persistence.
\newblock In \emph{Representations of Algebras and Related Structures},
EMS Series of Congress Reports, pages 77--150. European Mathematical
Society, 2023.
\newblock \doi{10.4171/ECR/19/4}.

\bibitem{ScoccolaEtAl2024}
L.~Scoccola, S.~Setlur, D.~Loiseaux, M.~Carrière, and S.~Oudot.
\newblock Differentiability and optimization of multiparameter persistent
homology.
\newblock In \emph{Proceedings of the 41st International Conference on
Machine Learning}, volume 235 of \emph{Proceedings of Machine Learning
Research}, pages 43986--44011, 2024.
\newblock \url{https://proceedings.mlr.press/v235/scoccola24a.html}.

\bibitem{BotnanOppermannOudot2025}
M.~B.~Botnan, S.~Oppermann, and S.~Oudot.
\newblock Signed barcodes for multi-parameter persistence via rank
decompositions and rank-exact resolutions.
\newblock \emph{Foundations of Computational Mathematics},
25:1815--1874, 2025.
\newblock \doi{10.1007/s10208-024-09672-9}.

\bibitem{BjerkevikBotnan2025}
H.~B.~Bjerkevik and M.~B.~Botnan.
\newblock Computing $p$-presentation distances is hard.
\newblock \emph{Discrete \& Computational Geometry}, 75(4):1414--1454,
2026; published online 2025.
\newblock \doi{10.1007/s00454-025-00744-3}.

\bibitem{LoiseauxSchreiber2024}
D.~Loiseaux and H.~Schreiber.
\newblock \texttt{multipers}: Multiparameter persistence for machine
learning.
\newblock \emph{Journal of Open Source Software}, 9(103):6773, 2024.
\newblock \doi{10.21105/joss.06773}.

\bibitem{deMouraUllrich2021}
L.~de Moura and S.~Ullrich.
\newblock The Lean~4 theorem prover and programming language.
\newblock In \emph{Automated Deduction---CADE 28}, pages 625--635.
Springer, 2021.
\newblock \doi{10.1007/978-3-030-79876-5_37}.

\bibitem{Mathlib2020}
The Mathlib Community.
\newblock The Lean mathematical library.
\newblock In \emph{Proceedings of the 9th ACM SIGPLAN International
Conference on Certified Programs and Proofs}, pages 367--381, 2020.
\newblock \doi{10.1145/3372885.3373824}.

\end{thebibliography}

\end{document}
